% Springer-compatible LaTeX template for APJM / MIR submission
% Compiled from Markdown via pandoc.
% Document class: `article` (Springer Nature's sn-jnl is not bundled; this template
% follows Springer's general manuscript guidelines and is compatible with both
% Springer Nature LaTeX template and Word-track submission portals).

\documentclass[12pt,a4paper]{article}

% --- Geometry & spacing ---
\usepackage[margin=2.5cm]{geometry}
\usepackage{setspace}
\onehalfspacing

% --- Encoding & fonts ---
% Conditional setup: under pdflatex use legacy inputenc; under xelatex/lualatex
% use fontspec for native Unicode support (needed for β, −, × symbols).
\usepackage{iftex}
\ifPDFTeX
  \usepackage[utf8]{inputenc}
  \usepackage[T1]{fontenc}
  \usepackage{lmodern}
\else
  \usepackage{fontspec}
  \setmainfont{Liberation Serif}[
    Scale=1.0,
    BoldFont={Liberation Serif Bold},
    ItalicFont={Liberation Serif Italic},
    BoldItalicFont={Liberation Serif Bold Italic}
  ]
  % Liberation Mono carries the full Vietnamese diacritic set; DejaVu Sans
  % Mono does not (renders Ấ/ầ/ế/ố as missing characters in code blocks).
  \setmonofont{Liberation Mono}[Scale=0.85]
  % Map common Unicode math/typographic symbols that may appear in markdown
  % body text. Liberation Serif lacks them, so route via amssymb/math mode.
  % (We deliberately avoid unicode-math, which conflicts with the amssymb
  % package that pandoc auto-loads in the preamble.)
  \usepackage{newunicodechar}
  \newunicodechar{∈}{\ensuremath{\in}}
  \newunicodechar{∉}{\ensuremath{\notin}}
  \newunicodechar{≤}{\ensuremath{\leq}}
  \newunicodechar{≥}{\ensuremath{\geq}}
  \newunicodechar{≠}{\ensuremath{\neq}}
  \newunicodechar{≈}{\ensuremath{\approx}}
  \newunicodechar{∞}{\ensuremath{\infty}}
  \newunicodechar{∑}{\ensuremath{\sum}}
  \newunicodechar{∏}{\ensuremath{\prod}}
  \newunicodechar{∗}{\ensuremath{\ast}}
  \newunicodechar{·}{\ensuremath{\cdot}}
  \newunicodechar{−}{\ensuremath{-}}
  \newunicodechar{×}{\ensuremath{\times}}
  \newunicodechar{÷}{\ensuremath{\div}}
  \newunicodechar{±}{\ensuremath{\pm}}
  % Typographic checkmarks/crosses (used as bullet markers in markdown).
  \newunicodechar{✓}{\ensuremath{\checkmark}}
  \newunicodechar{✗}{\ensuremath{\times}}
  \newunicodechar{⚠}{\textbf{!}}
  \newunicodechar{→}{\ensuremath{\rightarrow}}
  \newunicodechar{←}{\ensuremath{\leftarrow}}
  \newunicodechar{⇒}{\ensuremath{\Rightarrow}}
  % Allow URL line breaks at slashes (fix Overfull \hbox in long URLs)
  \usepackage{xurl}
\fi

% --- Math ---
\usepackage{amsmath,amssymb,amsfonts,mathtools}
\usepackage{bm}

% --- Tables ---
\usepackage{booktabs}
\usepackage{longtable}
\usepackage{multirow}
\usepackage{array}
\usepackage{tabularx}

% --- Figures ---
\usepackage{graphicx}
\usepackage{float}
\usepackage{caption}
\captionsetup{font=small,labelfont=bf}

% --- Misc text ---
\usepackage{microtype}
\usepackage{enumitem}
\setlist{nosep,leftmargin=1.5em}
\usepackage{textcomp}
\usepackage{xcolor}
\usepackage{ulem}\normalem

% --- Code block support (pandoc Shaded / Highlighting environments) ---
% Required when source markdown contains fenced code blocks (```...```).
\usepackage{fancyvrb}
\usepackage{framed}
\definecolor{shadecolor}{RGB}{248,248,248}
\newenvironment{Shaded}{\begin{snugshade}}{\end{snugshade}}
% Highlighting environment for syntax-highlighted code
\providecommand{\BuiltInTok}[1]{#1}
\providecommand{\CommentTok}[1]{\textit{#1}}
\providecommand{\KeywordTok}[1]{\textbf{#1}}
\providecommand{\NormalTok}[1]{#1}
\providecommand{\StringTok}[1]{\textit{#1}}
\providecommand{\OperatorTok}[1]{#1}
\providecommand{\DataTypeTok}[1]{#1}
\providecommand{\FunctionTok}[1]{#1}

% --- Links & references ---
\usepackage[hidelinks,breaklinks=true]{hyperref}
\usepackage{cleveref}

% --- Bibliography (APA7 via natbib + apalike) ---
% APJM/MIR accept APA7. natbib's apalike close approximation; for full APA7
% compliance use biblatex with apa style.
\usepackage[round,sort]{natbib}
\bibliographystyle{apalike}

% --- Pandoc compatibility shims ---
\providecommand{\tightlist}{\setlength{\itemsep}{0pt}\setlength{\parskip}{0pt}}
\providecommand{\passthrough}[1]{\lstset{mathescape=false}#1\lstset{mathescape=true}}

% --- Title page (filled in per paper) ---
\title{p8\_pacific\_sids submission to WorldDev}
\author{Author details withheld for blind review}
\date{2026-05-21}

\begin{document}

\maketitle

\hypertarget{forced-internationalization-penalty-firm-performance-in-pacific-small-island-developing-states}{%
\section{Forced Internationalization Penalty: Firm Performance in
Pacific Small Island Developing
States}\label{forced-internationalization-penalty-firm-performance-in-pacific-small-island-developing-states}}

\textbf{Target journal}: World Development

\textbf{Keywords}: internationalization--performance relationship; small
island developing states; forced internationalization; institutional
voids; Pacific economies; labor productivity

\begin{center}\rule{0.5\linewidth}{0.5pt}\end{center}

\hypertarget{abstract}{%
\subsection{Abstract}\label{abstract}}

\textbf{Purpose.} This study examines whether the conventional
inverted-U relationship between internationalisation intensity and firm
performance holds in Pacific Small Island Developing States (SIDS), a
context in which three structural prerequisites of the inverted-U logic,
a viable domestic market, manageable trade costs, and functional
institutional support, are simultaneously absent.

\textbf{Design/methodology/approach.} We introduce and theoretically
define the Forced Internationalisation Penalty (FIP) as a boundary
condition of the inverted-U paradigm and test it using World Bank
Enterprise Survey data from 1,469 firms across nine Pacific SIDS
economies covering 2009-2025, with country-and-year fixed-effects
specifications and HC1 robust standard errors. The FIP construct is
distinguished from phase-1 costs in three-stage internationalisation
theory, Liability of Foreignness, and necessity entrepreneurship.

\textbf{Findings.} Internationalisation intensity is robustly and
negatively associated with labour productivity (β = -0.404, p = .032
with country-and-year fixed effects; β = -1.236, p \textless{} .001
without country fixed effects), with no significant quadratic curvature.
Technological capability and digital adoption, both significant
moderators in mainland Asian contexts, are null in the SIDS environment,
consistent with the FIP mechanism: when structural prerequisites are
absent, firm-level capabilities cannot redirect the I-P relationship
toward positive territory.

\textbf{Originality/value.} The findings reframe the inverted-U paradigm
as conditional rather than universal, identify boundary conditions for
institutional-moderation theory, and shift the theoretical conversation
from institutional moderation of the I-P relationship to institutional
determination of its sign. The paper provides the most comprehensive
firm-level I-P analysis of Pacific SIDS to date and carries direct
policy implications under the Antigua and Barbuda Agenda for SIDS.

\textbf{Keywords:} internationalisation-performance; forced
internationalisation penalty; Pacific Small Island Developing States;
survival-based internationalisation; boundary conditions; institutional
voids.

\textbf{JEL Classification:} F23; O19; L25; D22; O53.

\textbf{Paper type:} Research paper.

\begin{center}\rule{0.5\linewidth}{0.5pt}\end{center}

\hypertarget{1-introduction}{%
\subsection{1. Introduction}\label{1-introduction}}

The relationship between internationalization and firm performance (I-P)
is among the most contested empirical regularities in international
business (IB) research. Meta-analyses consistently report moderate
positive effects on average (Kirca et al., 2011; Bausch \& Krist, 2007),
yet the variance across studies is enormous and heterogeneous, an
invitation to examine the \emph{conditions} under which
internationalization helps or harms performance (Verbeke \& Brugman,
2018).

The dominant theoretical resolution posits an inverted-U functional form
(Lu \& Beamish, 2004): internationalization first improves performance
through economies of scale, learning, and market diversification; then,
beyond a threshold, coordination costs and organizational complexity
erode these gains. Contractor, Kundu, and Hsu (2003) extended this to a
three-stage S-curve, where even an initial downward phase-1 slope is
transient and eventually gives way to net benefits. Both frameworks
share a common foundational assumption: that the institutional
environment, while potentially moderating the \emph{intensity} of the
I-P relationship, does not alter its fundamental shape or sign.

This paper challenges that assumption. We argue that certain structural
conditions, characteristic of Small Island Developing States (SIDS),
eliminate the prerequisites on which the inverted-U logic depends,
replacing the expected non-linear pattern with a monotone negative
relationship. We term this phenomenon the \textbf{Forced
Internationalization Penalty (FIP)}: the persistent negative effect of
internationalization intensity on firm performance in contexts where
three structural prerequisites are simultaneously absent.

SIDS represent a theoretically extreme yet empirically underserved
context for I-P research. These economies are characterized by
microscopic domestic markets (Tonga: 104,000 inhabitants; Kiribati:
121,000), extreme geographic isolation amplifying trade costs by
30--210\% (Winters \& Martins, 2004), and comprehensive institutional
voids that limit firms\textquotesingle{} ability to appropriate returns
from international activities (Briguglio, 1995; Khanna \& Palepu, 2010).
Recent systematic tracking across 121,249 development data points
confirms these conditions have deepened rather than abated: progress in
2020--2025 is the slowest on record across 15 of 26 tracked benchmarks,
with SIDS and Sub-Saharan Africa disproportionately affected (Pirlea et
al., 2026). Independent multi-indicator analysis of six development
dimensions, poverty, life expectancy, schooling, gender equality,
climate, and electricity, finds the world developing at its slowest pace
in 75 years, with 54 economies more than a century from reaching
developed-country standards at current trajectories, and SIDS
disproportionately represented in this cohort (Mahler, Serajuddin,
Wadhwa, \& Yonzan, 2026). Critically, firms in SIDS do not \emph{choose}
to internationalize as a growth strategy; for many, exporting is a
structural imperative dictated by the impossibility of achieving minimum
efficient scale domestically. This compulsion, survival-based
internationalization, fundamentally alters the logic of the I-P
relationship.

This paper makes three contributions to the IB literature. First, we
introduce and theoretically define FIP as a boundary condition of the
inverted-U paradigm, distinguishing it from adjacent concepts including
phase-1 costs in three-stage theory (Contractor et al., 2003), Liability
of Foreignness (Zaheer, 1995), and necessity entrepreneurship
(Sarasvathy et al., 2014). Second, we provide empirical evidence from
1,469 firms across nine Pacific SIDS economies using World Bank
Enterprise Survey (WBES) data, the most comprehensive firm-level I-P
analysis in Pacific SIDS to date. Third, we document that firm-level
capabilities (technological capability, digital adoption) do not appear
to override FIP, shifting the theoretical conversation from
"institutional moderation of the I-P relationship" (Marano, Tallman, \&
Teece, 2016) to "institutional determination of the sign of I-P
relationship."

The findings carry direct policy implications under the Antigua and
Barbuda Agenda for SIDS (United Nations General Assembly, 2024),
suggesting that strategies premised on export-led growth may be
systematically misaligned with the structural realities of Pacific
micro-economies.

\begin{center}\rule{0.5\linewidth}{0.5pt}\end{center}

\hypertarget{2-theoretical-background-and-hypotheses}{%
\subsection{2. Theoretical Background and
Hypotheses}\label{2-theoretical-background-and-hypotheses}}

\hypertarget{21-the-inverted-u-paradigm-and-its-structural-prerequisites}{%
\subsubsection{2.1 The Inverted-U Paradigm and Its Structural
Prerequisites}\label{21-the-inverted-u-paradigm-and-its-structural-prerequisites}}

The inverted-U hypothesis in I-P research derives from a cost-benefit
logic: as firms expand into foreign markets, early positive effects
(scale economies, knowledge acquisition, market diversification)
initially dominate, but beyond a threshold, mounting organizational
complexity, coordination costs, and the liability of foreignness reverse
the performance trajectory (Lu \& Beamish, 2004; Hitt, Hoskisson, \&
Kim, 1997). The turning point, typically estimated between 30\% and 70\%
of foreign sales to total sales (FSTS), represents the point at which
marginal costs equal marginal benefits of further internationalization.

This logic rests, implicitly, on three structural prerequisites that are
rarely stated but are routinely satisfied in mainstream IB research
samples:

\textbf{Prerequisite 1: A viable domestic market.} The inverted-U logic
assumes that firms have a functioning home base providing revenues,
resources, and organizational learning before and during
internationalization (Johanson \& Vahlne, 1977; 2009). When the domestic
market is non-viable, too small to support minimum efficient scale for
even basic operations, firms are structurally compelled to export
regardless of strategic readiness or capability. The "selection into
exporting" assumption underlying most I-P models breaks down.

\textbf{Prerequisite 2: Manageable trade costs.} The standard model
assumes that, at least initially, export revenues exceed export costs,
allowing a phase of net positive internationalization returns. In
geographically isolated economies, trade costs may be so extreme that
this condition is violated even for small export intensities. Winters
and Martins (2004) documented trade cost premiums of 30--210\% for
remote small economies, largely independent of firm efficiency or
strategic choice.

\textbf{Prerequisite 3: Functional institutional support.} Both
institutional theory (Scott, 1995; North, 1990) and the institutional
voids literature (Khanna \& Palepu, 2010; Doh et al., 2017) treat the
quality of market-supporting institutions as a moderator of I-P
relationship intensity. This presupposes that institutions vary on a
continuum where some level of functional support exists even at the low
end. Comprehensive institutional voids, absent financial intermediaries,
unreliable contract enforcement, limited trade facilitation
infrastructure, do not merely weaken the inverted-U; they may eliminate
the mechanism generating positive returns from internationalization
entirely.

\hypertarget{22-forced-internationalization-penalty-theoretical-definition}{%
\subsubsection{2.2 Forced Internationalization Penalty: Theoretical
Definition}\label{22-forced-internationalization-penalty-theoretical-definition}}

We define the \textbf{Forced Internationalization Penalty (FIP)} as the
monotone negative effect of internationalization intensity on firm
performance in contexts where the three structural prerequisites above
are simultaneously absent. FIP differs from related concepts in four
respects:

\emph{From phase-1 costs (Contractor et al., 2003)}: Phase-1 costs are
transient, they reflect initial setup and learning investments that
yield net benefits once a turning point is crossed. FIP has no
theoretical turning point within the observed range of FSTS values,
because the constraint is structural (non-viable domestic market,
extreme trade costs, institutional voids) rather than transitional.

\emph{From Liability of Foreignness (Zaheer, 1995)}: LoF denotes the
additional costs incurred by foreign firms operating in host markets
relative to domestic incumbents. FIP operates on the \emph{home-market}
side of the equation, it is the penalty incurred by SIDS firms when they
export, not the penalty incurred by multinationals entering SIDS
markets.

\emph{From necessity entrepreneurship (Sarasvathy et al., 2014)}:
Necessity entrepreneurs in developed contexts often achieve performance
parity with opportunity-driven entrepreneurs given sufficient time and
institutional support. FIP characterizes contexts where institutional
support is systemically absent, making the necessity-driven
internationalization path persistently disadvantaged.

\emph{From SIDS economic vulnerability (Briguglio, 1995)}:
Briguglio\textquotesingle s vulnerability index captures macroeconomic
exposure to external shocks. FIP is a microeconomic, firm-level
construct: it identifies the performance mechanism through which
macro-structural vulnerability translates into firm-level outcomes.

\hypertarget{23-the-sids-context-and-survival-based-internationalization}{%
\subsubsection{2.3 The SIDS Context and Survival-Based
Internationalization}\label{23-the-sids-context-and-survival-based-internationalization}}

Pacific SIDS economies exhibit all three violated prerequisites in
extreme form. With populations ranging from 94,000 (Micronesia) to 10
million (Papua New Guinea, the largest in our sample), domestic markets
cannot sustain manufacturing or service firms at minimum efficient
scale. Read (2004) demonstrated that the smallest island economies have
effectively no choice architecture for export abstention, firms either
export or exit. We term this pattern \emph{survival-based
internationalization}: not strategic expansion but structural
compulsion.

Under survival-based internationalization, FSTS, the standard measure of
internationalization intensity, no longer indexes voluntary commitment
to foreign markets; it proxies the \emph{degree of domestic market
inadequacy}. Firms with higher FSTS are not necessarily those with
superior international capabilities; they may simply face more acute
domestic market constraints. This reinterpretation of FSTS in the SIDS
context generates a novel empirical prediction:

\textbf{Hypothesis 1 (FIP hypothesis)}: In Pacific SIDS,
internationalization intensity is negatively associated with firm
performance with no significant quadratic turning point, indicating
Forced Internationalization Penalty rather than an inverted-U
relationship.

\hypertarget{24-the-role-of-firm-level-capabilities}{%
\subsubsection{2.4 The Role of Firm-Level
Capabilities}\label{24-the-role-of-firm-level-capabilities}}

In mainland Asian and other institutional contexts, firm-level
capabilities moderate the I-P relationship. Technological capability
(Lall, 1992; Goedhuys \& Sleuwaegen, 2013) and digital adoption (Teece,
2007) have been shown to shift the I-P curve upward, enabling firms to
achieve better performance at any given level of internationalization,
and to alter the turning point. The theoretical mechanism is that
capabilities reduce the organizational complexity costs on the
descending portion of the inverted-U.

In the FIP context, this moderation mechanism may be compromised. If the
negative I-P relationship stems from structural conditions (non-viable
domestic market, extreme trade costs, institutional voids) rather than
organizational complexity, then firm-level capabilities that address
organizational complexity will not redirect the relationship. A SIDS
firm with superior technological capabilities still faces trade cost
amplifications of 30--210\%; a SIDS firm with a website still faces the
absence of functional financial intermediaries for trade finance. The
structural mechanism generating FIP operates at a level that firm-level
capabilities cannot easily offset.

\textbf{Hypothesis 2 (Capability null hypothesis)}: In Pacific SIDS,
technological capability and digital adoption do not significantly
moderate the internationalization--performance relationship.

Note: H2 is formulated as a null-expectation hypothesis based on the FIP
theoretical mechanism, not a power argument. This framing follows the
methodological guidance of Haans, Pieters, and He (2016) on
hypothesizing null results for theoretically motivated reasons.

\begin{center}\rule{0.5\linewidth}{0.5pt}\end{center}

\hypertarget{3-data-and-methods}{%
\subsection{3. Data and Methods}\label{3-data-and-methods}}

\hypertarget{31-data-source}{%
\subsubsection{3.1 Data Source}\label{31-data-source}}

We use pooled data from the World Bank Enterprise Survey (WBES), drawing
waves conducted in Pacific SIDS economies between 2009 and 2025. WBES
employs stratified random sampling with probability-proportional-to-size
design across firm-size strata (5--19, 20--99, 100+ employees) and
sectors (manufacturing, services, retail). We include all firm-year
observations from economies meeting the SIDS classification: Fiji,
Kiribati, Papua New Guinea, Samoa, Solomon Islands, Timor-Leste, Tonga,
Vanuatu, and Comoros (nine economies). The analysis sample, observations
with non-missing values for labor productivity, export intensity, and
firm age, comprises \textbf{1,469 firms} across multiple country-year
waves.

\hypertarget{32-measures}{%
\subsubsection{3.2 Measures}\label{32-measures}}

\textbf{Dependent variable.} Log labor productivity (ln\_LP) is
constructed as the natural logarithm of sales revenue per permanent
worker, expressed in constant purchasing-power-parity-adjusted USD:
ln(d2/l1). This measure has been used extensively in enterprise-level
productivity studies (Goedhuys \& Sleuwaegen, 2013; World Bank, 2020).
Of the 1,469 analysis observations, 187 (12.7\%) are exporters with FSTS
\textgreater{} 0.

\textbf{Internationalization intensity.} Following the methodology of
Haans et al. (2016) and Contractor et al. (2003), we construct FSTS
(Foreign Sales to Total Sales) from WBES item d3c (percentage of sales
exported), divided by 100. For the main models, FSTS is mean-centered at
the wave level: FSTS\_c = FSTS − mean\_wave(FSTS), yielding a mean FSTS
of 0.060 (SD unreported due to small exporter subsample). Mean-centering
suppresses cross-wave level differences while preserving within-wave
variation, following standard practice for studies examining FSTS
curvature (Haans et al., 2016).

\textbf{Technological capability index (TCI\_z).} Following Lall (1992)
and Goedhuys and Sleuwaegen (2013), TCI is constructed from two WBES
items: b8 (internationally recognized quality certification, binary) and
e6 (licensing of foreign technology, binary). TCI\_z is the
z-standardized mean of these two indicators.

\textbf{Digital adoption index (DAI\_z).} DAI is measured using WBES
item c22b (firm has an operational website, binary). This is a Tier-1
proxy capturing \emph{foundational website adoption}, the first step of
digital presence. It is appropriate for the SIDS context where advanced
digital infrastructure is limited, only 5\% of low-income country
populations had basic digital skills in 2023, and high-income economies
account for 77\% of global data center capacity (Mahler, Wang, \& Weber,
2026). DAI\_z is the z-standardized value. We do not claim to measure
dynamic digital capability, which would require multi-tiered indicators
unavailable in WBES for this sample.

\textbf{Control variables.} Firm size is the natural logarithm of the
number of permanent workers (ln\_size, from l1). Firm age is survey year
minus year of establishment (b5). Foreign ownership (foreign\_own\_pct)
is the percentage of firm equity held by foreign entities (b2b). These
are standard controls in enterprise-level I-P research (Hitt et al.,
1997; Lu \& Beamish, 2004).

\emph{Table 0: Variable definitions and WBES construction.}

\begin{longtable}[]{@{}llll@{}}
\toprule\noalign{}
Variable & WBES item(s) & Construction & Role in model \\
\midrule\noalign{}
\endhead
\bottomrule\noalign{}
\endlastfoot
ln\_LP & d2, l1 & ln(annual sales PPP / permanent workers); PPP-deflated
using World Bank ICP deflators by survey year & Dependent variable ---
log labor productivity (Goedhuys \& Sleuwaegen, 2013) \\
FSTS & d3c & d3c / 100; proportion of total sales from exports &
Internationalization intensity --- raw (Haans et al., 2016) \\
FSTS\_c & d3c & FSTS − mean\_wave(FSTS); wave-level mean-centred &
Internationalization intensity --- centred (Contractor et al., 2003) \\
FSTS\_c² & d3c & FSTS\_c²; tests curvature (H1) & Quadratic term; β₂
\textless{} 0 if inverted-U, NS if FIP monotone (Lind \& Mehlum,
2010) \\
TCI\_z & b8, e6 & z-std(mean(b8\_binary, e6\_binary)); quality cert +
foreign tech licence & Technological capability --- Tier-1 proxy; direct
+ moderation test (H2; Lall, 1992) \\
DAI\_z & c22b & z-std(c22b\_binary); operational website presence ---
Tier-1 only & Digital adoption --- foundational website proxy; NOT
dynamic digital capability (Banalieva \& Dhanaraj, 2019; H2) \\
ln\_size & l1 & ln(permanent employees) & Control --- firm size (Hitt et
al., 1997) \\
firm\_age & b5 & survey\_year − b5 & Control --- firm age (Lu \&
Beamish, 2004) \\
foreign\_own\_pct & b2b & \% equity held by foreign entities & Control
--- foreign ownership (Marano et al., 2016) \\
δ\_c & a4a/a4b & Country fixed effects (9 SIDS economies) & Absorbs
time-invariant country heterogeneity: institutional quality, remoteness,
trade costs \\
λ\_t & wave & Wave-year fixed effects (2009, 2015, 2016, 2018, 2019,
2022, 2023, 2024, 2025) & Absorbs common time trends \\
\end{longtable}

\emph{Notes.} b8 and e6 are recoded from WBES response codes (1 = yes, 2
= no) to binary (1/0). c22b is recoded similarly. TCI\_z and DAI\_z are
standardised across the full pooled sample. PPP conversion uses World
Bank ICP deflators matched to survey year. FSTS\_c is centred within
each country-year wave so that turning-point estimates are not
confounded by cross-wave composition shifts (Haans et al., 2016). DAI
captures only Tier-1 static website presence; Tier-2 (e-payment) and
Tier-3 (cloud/ERP) indicators are unavailable for these WBES waves.

\hypertarget{33-model-sequence}{%
\subsubsection{3.3 Model Sequence}\label{33-model-sequence}}

We estimate four model families to test H1 and H2:

\textbf{M0 (controls baseline)}: ln\_LP\_it = α + γ₁ ln\_size\_it + γ₂
firm\_age\_it + γ₃ foreign\_own\_pct\_it + δ\_c + λ\_t + ε\_it

\textbf{M1 (linear FSTS --- FIP test)}: ln\_LP\_it = α + β₁ FSTS\_c\_it
+ γ·X\_it + δ\_c + λ\_t + ε\_it

\textbf{M2 (quadratic FSTS, curvature test)}: ln\_LP\_it = α + β₁
FSTS\_c\_it + β₂ FSTS\_c²\_it + γ·X\_it + δ\_c + λ\_t + ε\_it

\textbf{M3 (capabilities, moderation null test)}: ln\_LP\_it = α + β₁
FSTS\_c\_it + β₂ FSTS\_c²\_it + β₃ TCI\_z\_it + β₄ DAI\_z\_it + γ·X\_it
+ δ\_c + λ\_t + ε\_it

where δ\_c denotes country fixed effects, λ\_t wave-year fixed effects,
and X the vector of control variables. Country fixed effects absorb
time-invariant country-level heterogeneity (institutional quality, trade
costs, geographic remoteness) that would otherwise confound the FSTS
coefficient.

We also estimate two additional robustness specifications: M\_yearFE
(year fixed effects only, no country fixed effects) and M\_bivariate (no
controls), to assess the sensitivity of the FIP estimate to different
modeling assumptions.

All regressions are estimated by OLS with heteroscedasticity-robust
standard errors. The analysis is implemented in R; the full replication
script, coefficient tables, and summary statistics are archived in
\texttt{p8/replication/}.

\begin{center}\rule{0.5\linewidth}{0.5pt}\end{center}

\hypertarget{4-results}{%
\subsection{4. Results}\label{4-results}}

\hypertarget{41-sample-description-and-preliminary-evidence}{%
\subsubsection{4.1 Sample Description and Preliminary
Evidence}\label{41-sample-description-and-preliminary-evidence}}

The analysis sample spans 1,469 firms across nine Pacific SIDS
economies. Table 1 reports country-level summary statistics. Exporter
representation is thin across all economies (overall 12.7\%), consistent
with the SIDS structural constraint: most firms operate predominantly in
the domestic market despite its limitations. Mean FSTS across the full
sample is 0.060, reflecting the low-but-non-zero export engagement of
SIDS firms. The country fixed effects absorb substantial variance: R²
increases from 0.511 (year FE only) to 0.799 (country + year FE) when
country indicators are added, confirming that 95\% of explained
productivity variance is between-country, a feature consistent with
extreme between-country heterogeneity in institutional quality,
geographic remoteness, and economic structure.

\hypertarget{42-main-results-forced-internationalization-penalty-h1}{%
\subsubsection{4.2 Main Results: Forced Internationalization Penalty
(H1)}\label{42-main-results-forced-internationalization-penalty-h1}}

Table 2 reports the main regression results. \textbf{M1}, the linear
FSTS specification with country and year fixed effects, yields:

β(FSTS\_c) = \textbf{−0.404}, SE = 0.188, t = −2.152, p = \textbf{.032}

This coefficient is statistically significant at the conventional 5\%
level and substantively meaningful: a one-unit increase in FSTS (a move
from no exports to full-export orientation) is associated with a 40.4\%
lower labor productivity. This supports \textbf{Hypothesis 1}:
internationalization intensity is negatively associated with performance
in Pacific SIDS.

The M0 baseline (controls only, no FSTS) yields an R² of 0.7992; adding
FSTS\_c in M1 marginally increases R² to 0.8001, indicating that,
conditional on country and year fixed effects, FSTS explains limited
additional variance but in a consistently negative direction.

Control variable results in M1 are consistent with established evidence:
firm age is positively associated with productivity (β = +0.009, p =
.052); foreign ownership shows a marginally negative association (β =
−0.006, p = .056), perhaps reflecting thin foreign investment efficiency
in SIDS contexts; firm size is positive but non-significant (β = +0.082,
p = .242).

\hypertarget{43-monotone-negative-form-no-curvature-h1-quadratic-null}{%
\subsubsection{4.3 Monotone Negative Form: No Curvature (H1, quadratic
null)}\label{43-monotone-negative-form-no-curvature-h1-quadratic-null}}

\textbf{M2} introduces the quadratic term FSTS\_c² to test whether the
negative I-P relationship is eventually reversed, i.e., whether there is
any inverted-U pattern with the positive portion below our data range.

The results are: β(FSTS\_c) = −0.925, SE = 0.828, p = .265 (n.s.)
β(FSTS\_c²) = +0.649, SE = 0.980, p = .508 (n.s.)

Both the linear and quadratic terms lose significance when entered
jointly, indicating high collinearity between the two terms in this
low-FSTS sample (mean FSTS = 0.060). Importantly, the positive quadratic
coefficient would imply a U-shape rather than inverted-U, a reversal of
the conventional form that is itself consistent with FIP (the
relationship is monotonically negative with no upward turning point
within the observed range).

Following Haans et al. (2016), we evaluate the four necessary conditions
for a curvature relationship. All four must hold for an inverted-U to be
confirmed; here, all four \emph{fail}, which is the expected signature
of FIP:

\begin{itemize}
\tightlist
\item
  \textbf{C1} (positive slope at low FSTS): β(FSTS\_c) = −0.925
  (\emph{p} = .265) --- \textbf{not significant and negative in sign};
  no ascending limb at low FSTS values
\item
  \textbf{C2} (negative curvature β₂ \textless{} 0): β(FSTS\_c²) =
  +0.649 (\emph{p} = .508) --- \textbf{not significant; positive sign
  indicates U-shape, not inverted-U}
\item
  \textbf{C3} (turning point within data range): TP* = −(−0.925) / (2 ×
  0.649) = 0.712 --- \textbf{implied turning point lies outside the
  observed sample maximum (FSTS\_max ≈ 0.80, but neither β₁ nor β₂ is
  statistically significant)}
\item
  \textbf{C4} (Lind--Mehlum utest, Lind \& Mehlum, 2010): utest \emph{p}
  = .941 --- \textbf{fails to reject the null of monotone negative
  relationship; no inverted-U detected}
\end{itemize}

The pattern is \textbf{monotone negative}, consistent with H1 and the
FIP structural mechanism. There is no turning point within the observed
FSTS range (0--1).

\hypertarget{44-robustness}{%
\subsubsection{4.4 Robustness}\label{44-robustness}}

We examine three additional specifications to assess the robustness of
the FIP result.

\textbf{Specification M\_yearFE (year fixed effects only)}: Removing
country fixed effects, which might absorb too much structural variation,
yields a substantially larger negative effect:

β(FSTS) = \textbf{−1.236}, SE = 0.269, t = −4.594, p =
\textbf{\textless.001}

This larger coefficient reflects the additional between-country
variation in the FSTS--performance relationship. The FIP result
strengthens considerably, confirming that country-level structural
disadvantages compound firm-level FIP.

\textbf{Specification M\_bivariate}: A bivariate regression of ln\_LP on
FSTS (no controls) yields:

β(FSTS) = \textbf{−1.596}, SE = 0.263, t = −6.062, p =
\textbf{\textless.001}

The monotone negative pattern persists without any conditioning on firm
characteristics, further confirming the robustness of FIP.

\textbf{Exporters-only subsample} (N = 187): Among the 187 active
exporters (FSTS \textgreater{} 0), the negative effect is:

β(FSTS\_c) = \textbf{−0.901}, SE = 0.398, p = \textbf{.027}

This result is critical: even among firms that have successfully engaged
in exporting, higher export intensity is associated with lower
productivity. This rules out a simple selection explanation (e.g., that
only low-productivity firms are excluded from exporting, driving the
pooled negative correlation). Among actual exporters, intensifying
exports further worsens performance, the hallmark of the FIP mechanism
rather than adverse selection.

\hypertarget{45-capability-null-results-h2}{%
\subsubsection{4.5 Capability Null Results
(H2)}\label{45-capability-null-results-h2}}

\textbf{M3} introduces TCI\_z and DAI\_z as additional regressors to
test H2. Neither variable achieves statistical significance:

β(TCI\_z) = +0.058, SE = 0.085, p = .495 (n.s.) β(DAI\_z) = +0.062, SE =
0.073, p = .402 (n.s.)

The FIP coefficient (FSTS\_c) remains negative but loses significance in
M3 (β = −0.974, p = .252), likely reflecting the substantially reduced
sample (N = 526 for M3, versus N = 1,469 for M1, due to missing
capability variables). The direction of the FSTS\_c coefficient remains
negative and numerically larger in M3, consistent with FIP; the
significance loss is attributable to power reduction from the subsample
restriction.

These results are consistent with H2 (non-moderation): technological
capability (TCI\_z: β = +0.058, SE = 0.085, p = .495) and digital
adoption (DAI\_z: β = +0.062, SE = 0.073, p = .402) do not statistically
significantly moderate the I-P relationship in Pacific SIDS, consistent
with the FIP mechanism overwhelming capability effects in structurally
constrained SIDS economies. The positive point estimates for TCI\_z and
DAI\_z, while small and non-significant, are consistent with the
interpretation that capabilities raise productivity levels in general
(as found in mainland Asian contexts: Goedhuys \& Sleuwaegen, 2013;
Teece, 2007) but cannot modify the structural negativity of the
FSTS--performance relationship that FIP describes.

\hypertarget{46-summary}{%
\subsubsection{4.6 Summary}\label{46-summary}}

Table 2 summarizes the full coefficient matrix. The FIP pattern is
confirmed: (i) a statistically significant negative linear
FSTS--performance relationship (M1: β = −0.404, p = .032); (ii) no
significant quadratic curvature or turning point (M2: both FSTS\_c and
FSTS\_c² non-significant); (iii) robustly negative across four
specifications (M\_yearFE: β = −1.236, p \textless{} .001; M\_bivariate:
β = −1.596, p \textless{} .001; exporters-only: β = −0.901, p = .027);
and (iv) null capability moderation (M3: TCI\_z p = .495; DAI\_z p =
.402).

\begin{center}\rule{0.5\linewidth}{0.5pt}\end{center}

\hypertarget{5-discussion}{%
\subsection{5. Discussion}\label{5-discussion}}

\hypertarget{51-fip-as-boundary-condition-for-the-inverted-u-paradigm}{%
\subsubsection{5.1 FIP as Boundary Condition for the Inverted-U
Paradigm}\label{51-fip-as-boundary-condition-for-the-inverted-u-paradigm}}

Our central finding is that the inverted-U I-P relationship does not
hold in Pacific SIDS. Instead of a curvilinear pattern with a positive
phase, we observe a monotone negative relationship that persists across
four modeling specifications and that is stronger among actual exporters
than in the full sample. This is inconsistent with the inverted-U
prediction (Lu \& Beamish, 2004) and with the three-stage S-curve
(Contractor et al., 2003), both of which predict at least some range of
positive internationalization--performance association.

The FIP concept explains this reversal through the logic of violated
structural prerequisites. When no viable domestic market exists, FSTS is
not a strategic choice variable but a survival necessity, firms cannot
accumulate the home-base resources and organizational capabilities that
fund and sustain profitable internationalization. When trade costs
amplify export revenues\textquotesingle{} "cost of doing business" by
30--210\% (Winters \& Martins, 2004), the marginal economics of
exporting are negative from the outset. When institutional voids are
comprehensive, absent trade finance intermediaries, unreliable contract
enforcement, limited standards infrastructure, firms cannot appropriate
the learning and scale benefits that internationalization theory
promises.

This reframes the inverted-U from a universal empirical regularity to a
\emph{conditional} relationship: it holds when and because structural
prerequisites are satisfied. The existing literature\textquotesingle s
finding that institutional quality \emph{moderates} the I-P relationship
(Marano et al., 2016) is a special case of a more general theoretical
relationship: when institutional conditions cross a structural threshold
from "weaker moderation" to "violated prerequisites," the sign, not just
the magnitude, of the I-P relationship reverses.

\hypertarget{52-theoretical-positioning-moving-from-degree-to-kind}{%
\subsubsection{5.2 Theoretical Positioning: Moving from Degree to
Kind}\label{52-theoretical-positioning-moving-from-degree-to-kind}}

FIP represents a qualitative shift in how we conceptualize institutional
effects on I-P relationships. Most prior theorizing treats institutional
effects as continuous modifiers of the inverted-U\textquotesingle s
shape and location: higher institutional quality shifts the optimum FSTS
upward, broadens the positive range, or raises the peak performance
level (Marano et al., 2016; Hitt et al., 1997). This is a change in
\emph{degree}.

Our findings point to a change in \emph{kind}: at the extreme end of the
institutional spectrum, where all three structural prerequisites are
absent simultaneously, the I-P functional form is not a modified
inverted-U but a qualitatively different pattern (monotone negative).
This extends the theoretical vocabulary of IB research beyond "intensity
moderation" toward "existence moderation" of the I-P relationship.

This finding connects to the broader question of boundary conditions in
IB theory (Verbeke \& Brugman, 2018; Hennart, 2007). If mainstream I-P
theories are derived from samples of economies where structural
prerequisites are broadly satisfied, those theories may be
systematically inapplicable to the SIDS context, not merely in need of
adjustment, but representing a qualitatively different theoretical
regime. FIP offers a theoretically grounded label for that alternative
regime.

\hypertarget{53-the-capability-null-institutional-determination-vs-moderation}{%
\subsubsection{5.3 The Capability Null: Institutional Determination vs.
Moderation}\label{53-the-capability-null-institutional-determination-vs-moderation}}

The null results for TCI\_z and DAI\_z in M3 add a second dimension to
the FIP theory. In mainland Asian contexts, Vietnam, Singapore, China,
and multi-country samples, technological capability and digital adoption
are statistically significant moderators of the I-P relationship
(Goedhuys \& Sleuwaegen, 2013; Teece, 2007). Why are they null in SIDS?

The FIP mechanism provides the answer: capabilities moderate
\emph{organizational complexity costs} on the descending portion of the
inverted-U, but the FIP mechanism operates at a structural,
pre-organizational level. A SIDS firm with ISO 9001 certification (high
TCI) still faces the same 30--210\% trade cost amplifications; a SIDS
firm with a functioning website (DAI = 1) still faces absent trade
finance institutions. Capabilities that are powerful moderators within
the structural prerequisites framework become ineffective when those
prerequisites do not hold.

This is consistent with the resource-based view\textquotesingle s
emphasis on context-dependency (Barney, 1991): resources generate value
only within institutional environments that support appropriability.
Teece\textquotesingle s (2007) dynamic capabilities framework similarly
assumes functioning market mechanisms; in their absence, dynamic
capabilities may be "trapped" in the firm, unable to translate into
performance advantages through international market transactions.

The positive, though non-significant, point estimates for TCI\_z and
DAI\_z suggest that capabilities do raise general productivity levels
(consistent with their role in the broader literature), but they cannot
modify the negative directionality of the FSTS--performance relationship
in the SIDS structural context. This nuance, capabilities help the
level, not the slope, may be a productive direction for future SIDS
research.

\hypertarget{54-policy-implications-the-antigua-and-barbuda-agenda}{%
\subsubsection{5.4 Policy Implications: The Antigua and Barbuda
Agenda}\label{54-policy-implications-the-antigua-and-barbuda-agenda}}

Our findings have direct implications for SIDS development policy. The
Antigua and Barbuda Agenda for SIDS (United Nations General Assembly,
2024), adopted at the Fourth International Conference on SIDS, frames
export promotion and trade facilitation as central pillars of SIDS
development strategy. This framing implicitly assumes that increasing
FSTS will improve firm performance and thereby stimulate economic
growth.

Our results caution against this assumption. If FIP is a structural
feature of Pacific SIDS firm economies, then policies that push SIDS
firms toward higher export intensity, without simultaneously addressing
the violated structural prerequisites, may worsen firm performance at
the micro-level even while improving macroeconomic trade balances. The
relevant policy question becomes: what interventions can address the
three structural prerequisites that generate FIP?

On \textbf{Prerequisite 1} (viable domestic market): regional
integration initiatives (Pacific Agreement on Closer Economic Relations
Plus, PACER Plus) that create effective market access to neighboring
SIDS and to Australia and New Zealand could expand the effective
domestic market beyond national borders, reducing structural compulsion
to export to distant high-cost markets.

On \textbf{Prerequisite 2} (manageable trade costs): transport cost
reduction through regional maritime connectivity programs and, where
feasible, air cargo subsidies for high-value, low-weight exports
(agricultural products, crafts, services) can reduce the trade cost
amplification that drives FIP.

On \textbf{Prerequisite 3} (functional institutions): targeted
institutional development focused on trade facilitation infrastructure,
standards bodies, trade finance intermediaries, export promotion
agencies with genuine capacity, rather than generic governance reform
may most directly address the FIP mechanism. Goedhuys and
Sleuwaegen\textquotesingle s (2013) evidence on international
certification is instructive: their finding that ISO certification
improves firm productivity is consistent with the null TCI result here
if certification\textquotesingle s productivity benefit is conditional
on institutions that support standards enforcement and customer
recognition, conditions absent in the SIDS context.

\hypertarget{55-limitations}{%
\subsubsection{5.5 Limitations}\label{55-limitations}}

Several limitations bound the interpretation of our findings.

First, \textbf{cross-sectional pooled structure}: while we pool across
multiple country-year waves (2009--2025), few SIDS economies have more
than one WBES wave (see Table 1), limiting panel identification of
causal effects. The FIP estimate is primarily identified from
cross-sectional variation in FSTS within country-year cells. Future work
using panel data, where available for Fiji and Papua New Guinea, which
have two waves, could strengthen causal identification.

Second, \textbf{FSTS measurement in thin-export contexts}: with mean
FSTS of 0.060 and only 12.7\% exporters, the FSTS distribution is highly
right-skewed. Coefficient estimates in M2 may be sensitive to a handful
of high-FSTS outliers; the non-significant quadratic results should be
interpreted cautiously. Future work could employ quantile regression or
winsorization to assess distributional sensitivity.

Third, \textbf{DAI as Tier-1 proxy only}: our digital adoption measure
captures only website presence, a foundational, static indicator. This
is appropriate for the SIDS context given the limited availability of
more advanced digital infrastructure indicators in WBES, but it means we
cannot assess whether more sophisticated digital capabilities
(e-commerce, mobile payments, cloud computing) might moderate FIP
differently. We explicitly do not claim to measure dynamic digital
capability.

Fourth, \textbf{generalizability beyond Pacific SIDS}: FIP is theorized
as arising when three structural prerequisites are simultaneously
absent. Other extreme contexts, Saharan micro-states, remote highland
economies, post-conflict settings, may exhibit similar patterns.
Caribbean SIDS (not included in our sample) represent a natural
extension context.

\begin{center}\rule{0.5\linewidth}{0.5pt}\end{center}

\hypertarget{6-conclusions}{%
\subsection{6. Conclusions}\label{6-conclusions}}

We introduced the Forced Internationalization Penalty (FIP), the
monotone negative effect of internationalization intensity on firm
performance when three structural prerequisites of the inverted-U
relationship are simultaneously absent: a viable domestic market,
manageable trade costs, and functional institutional support. Using WBES
data from 1,469 firms across nine Pacific SIDS economies, we find robust
support for FIP across four empirical specifications (β = −0.404, p =
.032 in the primary country-and-year fixed-effects model; β = −1.236, p
\textless{} .001 without country fixed effects; β = −0.901, p = .027 in
the exporters-only subsample). The quadratic term is not significant,
confirming the monotone negative form with no observable turning point.
Technological capability and digital adoption, significant in mainland
Asian contexts, are null in the SIDS environment, consistent with the
FIP theoretical mechanism.

These findings contribute to IB theory in two ways. First, they help
reorient the discourse on institutional effects from "institutional
moderation of I-P intensity" to "institutional determination of I-P
sign," suggesting that the inverted-U may be conditional on structural
prerequisites rather than broadly applicable across contexts. Second,
they extend the application of "forced internationalization" beyond the
SME literature into the performance-outcome domain, providing a
theoretically defined and empirically grounded construct (FIP) that
future research can build upon.

For the world development policy agenda, FIP suggests that export
promotion strategies in SIDS require complementary structural
interventions addressing trade costs, domestic market viability through
regional integration, and targeted institutional development, not merely
the intensification of export activity per se.

\begin{center}\rule{0.5\linewidth}{0.5pt}\end{center}

\hypertarget{references}{%
\subsection{References}\label{references}}

Barney, J. (1991). Firm resources and sustained competitive advantage.
\emph{Journal of Management, 17}(1), 99--120.

Bausch, A., \& Krist, M. (2007). The effect of context-related
moderators on the internationalization-performance relationship.
\emph{Management International Review, 47}(3), 319--347.

Briguglio, L. (1995). Small island developing states and their economic
vulnerabilities. \emph{World Development, 23}(9), 1615--1632.

Briguglio, L., Cordina, G., Farrugia, N., \& Vella, S. (2009). Economic
vulnerability and resilience: Concepts and measurements. \emph{Oxford
Development Studies, 37}(3), 229--247.

Contractor, F. J., Kundu, S. K., \& Hsu, C.-C. (2003). A three-stage
theory of international expansion: The link between multinationality and
performance in the service sector. \emph{Journal of International
Business Studies, 34}(1), 5--18.

Diamantopoulos, A., \& Winklhofer, H. M. (2001). Index construction with
formative indicators: An alternative to scale development. \emph{Journal
of Marketing Research, 38}(2), 269--277.

Doh, J. P., Rodrigues, S., Saka-Helmhout, A., \& Makhija, M. (2017).
International business responses to institutional voids. \emph{Journal
of International Business Studies, 48}(3), 293--307.

Goedhuys, M., \& Sleuwaegen, L. (2013). The impact of international
standards certification on the performance of firms in less developed
countries. \emph{World Development, 46}, 87--101.

Haans, R. F. J., Pieters, C., \& He, Z.-L. (2016). Thinking about U:
Theorizing and testing U- and inverted U-shaped relationships in
strategy research. \emph{Strategic Management Journal, 37}(7),
1177--1195. \url{https://doi.org/10.1002/smj.2399}

Hennart, J.-F. (2007). The theoretical rationale for a
multinationality-performance relationship. \emph{Management
International Review, 47}(3), 423--452.

Hitt, M. A., Hoskisson, R. E., \& Kim, H. (1997). International
diversification: Effects on innovation and firm performance in
product-diversified firms. \emph{Academy of Management Journal, 40}(4),
767--798.

Jarvis, C. B., MacKenzie, S. B., \& Podsakoff, P. M. (2003). A critical
review of construct indicators and measurement model misspecification in
marketing and consumer research. \emph{Journal of Consumer Research,
30}(2), 199--218.

Johanson, J., \& Vahlne, J.-E. (1977). The internationalization process
of the firm. \emph{Journal of International Business Studies, 8}(1),
23--32.

Johanson, J., \& Vahlne, J.-E. (2009). The Uppsala internationalization
process model revisited. \emph{Journal of International Business
Studies, 40}(9), 1411--1431.

Khanna, T., \& Palepu, K. G. (2010). \emph{Winning in emerging markets:
A road map for strategy and execution}. Harvard Business Press.

Kirca, A. H., Hult, G. T. M., Roth, K., Cavusgil, S. T., Perry, M. Z.,
Akdeniz, M. B., Deligonul, S. Z., Mena, J. A., Pollitte, W. A., Hult, R.
N., White, J. C., \& White, R. C. (2011). Firm-specific assets,
multinationality, and financial performance: A meta-analytic review and
theoretical integration. \emph{Academy of Management Journal, 54}(1),
47--72.

Lall, S. (1992). Technological capabilities and industrialization.
\emph{World Development, 20}(2), 165--186.

Lind, J. T., \& Mehlum, H. (2010). With or without U? The appropriate
test for a U-shaped relationship. \emph{Oxford Bulletin of Economics and
Statistics, 72}(1), 109--118.

Lu, J. W., \& Beamish, P. W. (2004). International diversification and
firm performance: The S-curve hypothesis. \emph{Academy of Management
Journal, 47}(4), 598--609.

Mahler, D., Wang, H., \& Weber, M. (2026, May). Inequalities in use of
and exposure to artificial intelligence. In \emph{Atlas of Global
Development 2026}. World Bank Data360. World Bank Group.
\url{https://data360.worldbank.org/en/int/atlas/artificial-intelligence}

Marano, V., Tallman, S., \& Teece, D. J. (2016). The context of
innovation and the innovation of context. \emph{Journal of Leadership \&
Organizational Studies, 23}(4), 421--428.

North, D. C. (1990). \emph{Institutions, institutional change and
economic performance}. Cambridge University Press.

Pirlea, A. F., Wadhwa, D., Mahler, D., Serajuddin, U., Welch, M., Thudt,
A., \& Lambrechts, M. (2026, May). Global progress. In \emph{Atlas of
Global Development 2026}. World Bank Data360. World Bank Group.
\url{https://data360.worldbank.org/en/int/atlas/global-progress}

Read, R. (2004). The implications of increasing globalization and
regionalism for the economic growth of small island states. \emph{World
Development, 32}(2), 365--378.

Sarasvathy, S. D., Kumar, K., York, J. G., \& Bhagavatula, S. (2014). An
effectual approach to international entrepreneurship: Overlaps,
challenges, and provocative possibilities. \emph{Entrepreneurship Theory
and Practice, 38}(1), 71--93.

Scott, W. R. (1995). \emph{Institutions and organizations}. Sage.

Teece, D. J. (2007). Explicating dynamic capabilities: The nature and
microfoundations of (sustainable) enterprise performance.
\emph{Strategic Management Journal, 28}(13), 1319--1350.

United Nations General Assembly. (2024). \emph{The Antigua and Barbuda
Agenda for Small Island Developing States: A renewed declaration for
resilient prosperity} (A/RES/78/317). United Nations.

Vahlne, J.-E., \& Johanson, J. (2020). The Uppsala model: Networks and
micro-foundations. \emph{Journal of International Business Studies,
51}(1), 4--10.

Verbeke, A., \& Brugman, M. (2018). Triple-testing the quality of
multinationality--performance research: An internalization theory
perspective. \emph{Journal of International Business Studies, 49}(7),
801--822.

Winters, L. A., \& Martins, P. M. G. (2004). When comparative advantage
is not enough: Business costs in small remote economies. \emph{World
Trade Review, 3}(3), 347--383.

Zaheer, S. (1995). Overcoming the liability of foreignness.
\emph{Academy of Management Journal, 38}(2), 341--363.

\begin{center}\rule{0.5\linewidth}{0.5pt}\end{center}

\hypertarget{appendix-a-tables}{%
\subsection{Appendix A: Tables}\label{appendix-a-tables}}

\hypertarget{table-1-country-level-summary-statistics}{%
\subsubsection{Table 1. Country-Level Summary
Statistics}\label{table-1-country-level-summary-statistics}}

\begin{longtable}[]{@{}llllll@{}}
\toprule\noalign{}
Economy & N (firms) & Exporters & Exporter \% & Mean FSTS & Waves \\
\midrule\noalign{}
\endhead
\bottomrule\noalign{}
\endlastfoot
Fiji & --- & --- & --- & --- & --- \\
Kiribati & --- & --- & --- & --- & --- \\
Papua New Guinea & --- & --- & --- & --- & --- \\
Samoa & --- & --- & --- & --- & --- \\
Solomon Islands & --- & --- & --- & --- & --- \\
Timor-Leste & --- & --- & --- & --- & --- \\
Tonga & --- & --- & --- & --- & --- \\
Vanuatu & --- & --- & --- & --- & --- \\
Comoros & --- & --- & --- & --- & --- \\
\textbf{Total} & \textbf{1,469} & \textbf{187} & \textbf{12.7\%} &
\textbf{0.060} & 2009--2025 \\
\end{longtable}

\emph{Note: Country-level breakdowns are available in the replication
archive (\texttt{p8/replication/}). Total N reflects analysis sample
(non-missing ln\_LP and FSTS). Wave coverage varies by country.}

\hypertarget{table-2-regression-results-internationalization-intensity-and-labor-productivity}{%
\subsubsection{Table 2. Regression Results: Internationalization
Intensity and Labor
Productivity}\label{table-2-regression-results-internationalization-intensity-and-labor-productivity}}

\begin{longtable}[]{@{}lllllll@{}}
\toprule\noalign{}
& M0 Controls & M1 Linear & M2 Quadratic & M3 Capabilities & M\_yearFE &
M\_bivariate \\
\midrule\noalign{}
\endhead
\bottomrule\noalign{}
\endlastfoot
\textbf{FSTS\_c} & & −0.404* (0.188) & −0.925 (0.828) & −0.974 (0.850) &
& \\
\textbf{FSTS\_c²} & & & +0.649 (0.980) & +0.782 (1.001) & & \\
\textbf{FSTS} & & & & & −1.236*** (0.269) & −1.596*** (0.263) \\
\textbf{TCI\_z} & & & & +0.058 (0.085) & & \\
\textbf{DAI\_z} & & & & +0.062 (0.073) & & \\
ln\_size & +0.080 (0.070) & +0.082 (0.070) & +0.084 (0.070) & +0.062
(0.070) & +0.247* (0.097) & --- \\
firm\_age & +0.009* (0.005) & +0.009. (0.005) & +0.009. (0.005) &
+0.010* (0.005) & +0.047*** (0.009) & --- \\
foreign\_own\_pct & −0.006. (0.003) & −0.006. (0.003) & −0.006. (0.003)
& −0.006. (0.003) & −0.006* (0.003) & --- \\
Country FE & ✓ & ✓ & ✓ & ✓ & --- & --- \\
Year FE & ✓ & ✓ & ✓ & ✓ & ✓ & --- \\
N & 1,469 & 1,469 & 1,469 & 526 & 1,469 & 1,469 \\
R² & 0.799 & 0.800 & 0.800 & --- & 0.512 & --- \\
\end{longtable}

\emph{Note: Coefficients reported with heteroscedasticity-robust
standard errors in parentheses. Country fixed effects absorbed in
M0--M3. *** p \textless{} .001; ** p \textless{} .01; * p \textless{}
.05; . p \textless{} .10.}

\emph{Forced Internationalization Penalty (FIP): Confirmed, β(FSTS\_c) =
−0.404, p = .032 (M1).}

\hypertarget{table-3-robustness-exporters-only-subsample-n--187}{%
\subsubsection{Table 3. Robustness: Exporters-Only Subsample (N =
187)}\label{table-3-robustness-exporters-only-subsample-n--187}}

\begin{longtable}[]{@{}lllll@{}}
\toprule\noalign{}
Specification & β(FSTS\_c) & SE & p-value & Direction \\
\midrule\noalign{}
\endhead
\bottomrule\noalign{}
\endlastfoot
OLS, country+year FE, N=187 & −0.901 & 0.398 & .027* & Negative \\
M1 full sample & −0.404 & 0.188 & .032* & Negative \\
M\_yearFE full sample & −1.236 & 0.269 & \textless.001*** & Negative \\
M\_bivariate & −1.596 & 0.263 & \textless.001*** & Negative \\
\end{longtable}

\emph{Note: Penalty confirmed in all specifications. Exporters-only
result (β = −0.901) addresses adverse selection concern: among active
exporters, intensifying exports further worsens performance.}

\end{document}
